% !TeX root = ../papers/nc-p3.tex
\section{階数十の四次元分岐中心に続くねじれ三次曲線}
\label{sec:rank-ten-cubic-after-center}
\begin{theorem}\label{thm:rank-ten-cubic-after-center}
反対階数10の標準正則窓で曲面射影がなく、中央支持4の順中心 $p=x_0$ に始まる
列の第二曲線 $x_2$ がねじれ三次曲線とする。第三曲線の次数は4以下であり、第二曲線上に順中心がある。
第三曲線が非平面的なら第四曲線も動き、$p$ を通る平面に入り、
\[
             2\le d_3-1\le d_4\le3
\]
である。さらに第三曲線は滑らかな二次曲面の二つの母線族に対する射影次数が
$1,d_3-1$ の有理曲線である。逆方向でも成立する。
特に中央支持4の中心をどちらかの端に持つ次数列 $(1,3,3,1)$ で、
内側二本がともにねじれ三次曲線であるものはない。
\end{theorem}
\begin{proof}[二次の二項行列と一次の関係]
第一直線と第二曲線のパラメータを $s,t$ とする。$t=(t_0,t_1)$ とし
\[
              w(-,x_1(s),x_2(t),-)=[s,t]A(t).
\]
位置0の行支持は $U=\ker p$、位置3の列の変数空間は四次元 $V_3^*$ とする。
$A(t)$ の全列が張る $D\subset U\otimes k[t_0,t_1]_2$ は四次元で、係数の $U$ 支持は全三次元である。
実際、第二曲線が位置2を張るから、第一直線の一般点における中央支持と逆支持の全てがここに現れる。
逆支持は3、中央支持は4である。後者は定理\ref{thm:branch-binary-center-pairing}より、
この直線の中央支持が3ならその一意な前点の中心が中央支持2になることによる。
$w$ の反対切断で位置1を第一直線に制限するとその像は $k[t_0,t_1]_1D$ である。
これは反対側の $p$ での四次元像を消すので次元6以下である。

$S_j=k[t_0,t_1]_j$ とし $D_-=\{L\in U\otimes S_1:S_1L\subset D\}$、
$D_{--}=\{u\in U:S_1u\subset D_-\}$ と置く。
$t_0 f=t_1 g$ の関係を割ると $\dim S_1D=8-\dim D_-$、
$\dim S_1D_-=2\dim D_- -\dim D_{--}$ を得る。
前者から $m=\dim D_-\ge2$、後者から $2m-\dim D_{--}\le4$ である。
$2\dim D_{--}\le m$ もあるので $m=2$ である。
$D_{--}=0$ なら二つの一次列 $L_1,L_2$ に対し $D=S_1L_1+S_1L_2$、
$D_{--}=ku$ なら $D=S_2 u+kB(t)$ である。
両場合とも一般の値 $P(t)=A(t)V_3^*$ は二次元である。
最初の場合、二つの一次列が常に比例するなら、独立性からともに固定ベクトルの一次式倍となり
$D_{--}=0$ に反する。第二の場合は $U$ の完全支持による。

$R_0$ の第二位置を第一直線に制限する写像を考える。像は
$([-,t]\otimes A(t)z)_{t,z}$ が張る空間である。実際 $R_0$ は反対辺の二変数の
全縮約の像であり、$x_2(t)$ が全空間を張る。制限写像には非零核がある。
さもなくば $z\in\ker A(t)$ の全てが $w(-,-,x_2(t),z)=0$ となり、
動く一般点 $x_2(t)$ が順中心になってしまう。従って六次元の $S_1\otimes U$ における
像の消去式があり、一次行 $\ell(t)$ が全 $P(t)$ を消す。
$U$ の完全支持により定数の消去式はないので $\ell$ の二係数は独立である。
一般の $P(t)$ は $\ker\ell(t)$ に等しく、その全てが固定直線 $ku$ を含む。
一定基底を取り $P(t)=\langle u,v(t)\rangle$、$v(t)=t_0 v_0+t_1 v_1$ と書けば
$u,v_0,v_1$ は独立である。

これで次の二つの形で尽きる。$D_{--}=0$ なら、$L_i$ の $ku$ を法とする像は
それぞれ $v(t)$ の定数倍である。基底を取り直し一方を $n(t)u$、他方を $v(t)$ とできる
（後者の $u$ 方向は $v_0,v_1$ に吸収する）。$D_{--}\ne0$ の場合は $S_2u$ の
$u$ がまさに全 $P(t)$ に共通する直線を張り、$B(t)$ の剰余は $n(t)v(t)$、
$u$ 方向は他の三列で消せる。$n$ は両方で非零一次式である。位置3の列基底の選択により
\[
 \begin{array}{ll}
 \mathrm{I}: & A(t)=(t_0 n(t)u,\ t_1 n(t)u,\ t_0 v(t),\ t_1 v(t)),\\
 \mathrm{II}:& A(t)=(t_0^2 u,\ t_0t_1 u,\ t_1^2u,\ n(t)v(t)).
 \end{array}
\]
一般の $P(t)$ の消去行は上記の原始一次行の多項式倍なので、
第一直線への $R_0$ の制限の余次元はちょうど1、核も一次元である。
$n(t)=0$ では $A(t)$ がちょうど階数1になり、$w(-,-,x_2(t),-)$ の像の制限も一次元。
従って元の像は二次元以下である。二次元の下界より、ここが第二曲線上の順中心となる。
またそこでの $A(t)$ 自身は $R_3$ の純関係である。
\end{proof}
\begin{proof}[後続の曲線の次数]
第三点 $x_3(t)$ は $A(t)$ の核に入る。第三点が一定なら二直線列以下となり、
定理\ref{thm:two-line-full-middle-restriction}に反する。
II型なら $x_3$ は最後の座標の零平面に入り、平面曲線の次数上界3を使える
（系\ref{cor:rank-nine-nonlinear-corridors}の末尾の階数によらない証明）。
I型では、一般の核は二行をそれぞれ $(t_0,t_1)$ が消す二行二列行列であるから
\[
         x_3(t)=(-t_1 f(t),t_0 f(t),-t_1 g(t),t_0 g(t)).
\]
ここで $f,g$ は同次数で基点のない二項式で、次数は $d_3-1$。
射影直線から二次曲面の各因子への射として読むことで、特殊な $t$ にもこの表示が延びる。
第三曲線が平面なら既に十分である。非平面なら $f,g$ は独立であり、
$n=t_0$ とできる。第四点 $x_4(t)$ での三つの評価を $u(x_4),v_0(x_4),v_1(x_4)$ とする。
$x_3(t),x_4(t)$ は辺なので、第二位置には独立なパラメータ $y$ を代入できる。
$A(y)x_3(t)$ を $x_4(t)$ が消す式から $[t,y]$ を除くと
\[
 y_0\big(f(t)u(x_4)+g(t)v_0(x_4)\big)+y_1 g(t)v_1(x_4)=0
\]
を得る。従って $x_4$ は $v_1=0$ の平面に入り、$p$ からの射影は
$[u(x_4):v_0(x_4)]=[g:-f]$ である。
$x_4$ が恒等的に $p$ となることは、順中心 $p$ の前点が一意なので不可能。
射影が非定数なので第四曲線は動き、射影で次数は増えないから所望の不等式となる。
\end{proof}
